\documentclass[11pt]{article}
\usepackage[margin=1in]{geometry}
\usepackage{amsmath,amssymb,amsthm}
\usepackage{microtype}
\usepackage{parskip}
\usepackage[T1]{fontenc}
\usepackage{lmodern}

\begin{document}

\begin{center}
  {\large\bfseries RetroDiff: Deriving the Retrieval-Augmented Training Objective}
\end{center}

\medskip

We derive the training objective for \textbf{RetroDiff}, a retrieval-augmented diffusion
model. Let $x$ denote a clean data sample, $x_t$ its noised counterpart at step $t$,
$q$ the forward noising distribution, $p_\theta$ the learned reverse distribution with
parameters $\theta$, and $T$ the total number of diffusion steps.

\paragraph{Step 1: Variational lower bound.}
Because the marginal likelihood $p_\theta(x)$ is intractable, we lower-bound it via
Jensen's inequality applied to the expectation over the forward process
$q(x_{1:T}\mid x)$:

\begin{align}
  \log p_\theta(x)
  \;\geq\;
  \mathbb{E}_{q(x_{1:T}|x)}\!\left[\log\frac{p_\theta(x_{0:T})}{q(x_{1:T}|x)}\right]
  \;=\;
  -\mathcal{L}_{\mathrm{ELBO}}
  \label{eq:elbo}
\end{align}

Maximising $\log p_\theta(x)$ is therefore equivalent to minimising
$\mathcal{L}_{\mathrm{ELBO}}$.

\paragraph{Step 2: Decomposition into tractable terms.}
Expanding the joint $p_\theta(x_{0:T})=p(x_T)\prod_t p_\theta(x_{t-1}\mid x_t)$
and the forward posteriors $q(x_{t-1}\mid x_t,x)$, which are Gaussian in closed form,
yields the canonical decomposition:

\begin{align}
  \mathcal{L}_{\mathrm{ELBO}}
  \;=\;
  \underbrace{\mathbb{E}_q\!\left[-\log p_\theta(x_0\mid x_1)\right]}_{\text{reconstruction}}
  \;+\;
  \sum_{t=2}^{T}
  \underbrace{
    \mathrm{KL}\!\left(q(x_{t-1}\mid x_t,x)\,\|\,p_\theta(x_{t-1}\mid x_t)\right)
  }_{\text{step-}t\text{ denoising}}
  \label{eq:decomp}
\end{align}

The reconstruction term trains $p_\theta$ to recover $x_0$ from a single noisy
observation $x_1$; each KL term penalises deviation of the learned reverse step from
the tractable forward posterior.

\paragraph{Step 3: Retrieval-augmented objective (RetroDiff).}
Equations~\eqref{eq:elbo}--\eqref{eq:decomp} treat samples in isolation.
RetroDiff's key contribution is to condition the per-sample ELBO on a set
$R(x)$ of $k$ retrieved exemplars drawn from the training corpus.  Averaging over
the data distribution gives the final objective:

\begin{align}
  \mathcal{L}_{\mathrm{RetroDiff}}(\theta)
  \;=\;
  \mathbb{E}_{x\sim p_{\mathrm{data}}}\!\left[
    \mathcal{L}_{\mathrm{ELBO}}\!\left(x\;\big|\;R(x)\right)
  \right]
  \label{eq:retro}
\end{align}

By conditioning on $R(x)$, the model gains access to high-quality denoising hints at
every diffusion step, reducing the effective complexity of each KL term in
\eqref{eq:decomp} without modifying the model architecture.

\end{document}
